% ----------------------------------------------------------
% Introdução
% ----------------------------------------------------------
\chapter{INTRODUÇÃO}

A indústria de jogos digitais apresentou um crescimento expressivo ao longo dos últimos anos, ampliando sua relevância econômica, seu faturamento e sua base de usuários \cite{mbakirtzis2023analise}. Nesse cenário, as comunidades de jogadores passaram a assumir uma participação mais ativa, deixando de atuar apenas como consumidoras para também produzir e compartilhar conteúdos relacionados aos jogos \cite{soares2019jogos}.

\citeonline{fortim2022pesquisa} aponta que as desenvolvedoras enfrentam dificuldades para identificar o perfil dos jogadores, definir o público-alvo, fidelizá-lo e estabelecer estratégias de marketing adequadas. Nesse contexto, a aplicação de instrumentos de avaliação da experiência do jogador pode contribuir para reunir percepções sobre pontos fortes e fracos dos jogos, problemas técnicos, dificuldades de interação e aspectos positivos da experiência \cite{pinto2024avaliacao}.

A utilização dessas informações pode contribuir para a identificação de aspectos técnicos, pontos fortes, dificuldades de interação e elementos que podem ser aprimorados no jogo. Dessa forma, essas percepções podem auxiliar a tomada de decisão durante o desenvolvimento e aproximar o produto das expectativas do público.

Entretanto, o aproveitamento desse potencial pode ser dificultado pela diversidade de interações presentes em diferentes ambientes digitais. Essa variedade de finalidades pode tornar mais complexa a identificação de contribuições diretamente úteis ao desenvolvimento, especialmente diante das dificuldades das empresas em compreender seu público e organizar informações relacionadas aos usuários \cite{fortim2022pesquisa}.

Como produto deste estudo, será desenvolvido o protótipo de uma plataforma web voltada à organização de sugestões para jogos digitais, planejado para mediar a interação entre desenvolvedores e jogadores de diferentes nichos, promovendo o compartilhamento de ideias, sugestões e \textit{feedbacks} relacionados aos jogos. Adicionalmente, o sistema busca reunir e organizar informações que possam auxiliar empresas e desenvolvedores independentes na análise de características e aspectos de seus produtos, contribuindo para decisões mais fundamentadas durante o processo de desenvolvimento.
